% !TEX root = ../main.tex

Ho scritto tante cose in questa tesi, ma niente è stato più difficile dei ringraziamenti. All'inizio non volevo nemmeno scriverli: un po' perché non sono bravo con le parole in queste occasioni, un po' per evitare il momento sentimentale, ma soprattutto perché mi imbarazzano da morire. Parlo molto, lo sapete tutti, ma non sono fatto per queste cose. Nonostante tutto, mi sono convinto a dedicare un po' di spazio a chi mi ha accompagnato in questo percorso e mi ha aiutato a crescere.

Ai miei genitori, \textbf{mamma} e \textbf{papà}: se c'è qualcuno più felice di me per questo traguardo sono loro, per motivazioni diverse ma ugualmente importanti.

\textbf{Mamma}, mi hai sopportato tanto e non è stato per niente facile. In questi anni mi hai accompagnato ovunque e hai svuotato molti borsoni; sebbene la cucina chiudesse alle 20:00, hai fatto gli straordinari per farmi mangiare dopo le partite e i miei allenamenti assurdi. Ti ho fatto mille storie sulla dieta e su cosa mangiare e tu, nel bene o nel male, hai cercato di accontentarmi anche quando non eri d'accordo. Ogni giorno me ne uscivo con una nuova, ma se non altro non ti ho fatta annoiare. A volte ti ho dato pensieri e preoccupazioni, ma sei quella che mi conosce più a fondo. Solo tu riesci a capire cosa provo in ogni momento e intuisci subito se qualcosa non va. Mi vuoi bene più di ogni altra cosa e, quando sono triste, sei la prima a rimanerci male, assorbendo le mie emozioni. L'unica cosa che vuoi è vedermi felice e realizzato. Mi hai sempre incoraggiato a sorridere, consigliandomi spesso di staccare la spina da ciò che mi faceva stare male. Mi hai ricordato che nella vita c'è tanto altro e di non buttarmi mai giù, perché in cuor tuo mi hai sempre definito una \textit{“potenza”}. Non so quante volte tu mi abbia detto che mi seguirai ovunque andrò, ma dentro di me so che sei sempre al mio fianco. Adesso che sono un po' più libero e sta arrivando l'estate, possiamo tornare a prenderci un gelato appena torni da scuola, come facevamo ai tempi del liceo, e ti prometto che cercherò di sbattere meno il pugno sul tavolo. Ti voglio bene mamma, grazie di tutto.

\textbf{Papà}, io e te siamo molto diversi, ma anche tanto simili. Io odio fare la spesa e svegliarmi presto al mattino, mentre tu sei il primo ad alzarti e a pensare a cosa comprare per riempire dispensa e frigo. Spesso ti lamenti del fatto che io sia disordinato e smemorato; dici che ho preso da mamma, anche se lei sostiene il contrario. Hai un gusto discutibile nel vestire, ma ogni tanto cerchi di migliorare il tuo guardaroba trafugando nel mio. Nonostante tutto, sei sempre stato un esempio e, a volte senza nemmeno rendermene conto, prendo spunto dalle tue azioni e dai tuoi comportamenti per cercare la soluzione giusta ai problemi più difficili. In questi mesi ti ho costretto ad ascoltarmi mentre ripetevo per gli esami, e anche stavolta ti ho fatto leggere la tesi e sorbire la presentazione, nonostante tornassi stanco dal lavoro. Anche tu mi hai sempre spronato e incoraggiato, ricordandomi che questa è solo una fase della vita e che le soddisfazioni arriveranno poi. Speri sempre che l'unica persona al mondo a superarti sia io, e che possa raggiungere traguardi sempre più alti. Anche se non ci apriamo molto, ti voglio bene papà, grazie di tutto.

\textbf{Carry}, quante volte hai dormito mentre lavoravo e quante volte ti sei piazzata davanti allo schermo durante una videochiamata? Qualcuno potrebbe pensare che sia pigrizia o voglia di stare al centro dell'attenzione, ma in realtà stavi solo cercando di spiegarmi che il riposo è importante e volevi darmi supporto morale durante i colloqui. Ti definiscono una gattina antipatica e per niente affettuosa, ma ogni sera ti accoccoli sulle mie gambe, soprattutto quando sono giù di corda. Grazie anche a te.

Non potevo chiedere una famiglia migliore.

Oltre ai miei familiari, ci sono tante altre persone che meritano un posto in questi ringraziamenti. La prima è la ragione per cui ho deciso di scriverli (giuro, senza alcuna pressione).

\textbf{Giusy}, non sono mai andato d'accordo con le insegnanti e non ho mai avuto una grande opinione delle maestre. E infatti, non a caso, ho scelto un'insegnante, per di più una maestra, come fidanzata. Sarà deformazione professionale, ma mi hai insegnato tantissime cose: grazie a te ho scoperto che per una vita intera ho sbagliato a chiamare “asilo” ed “elementari” la scuola dell'infanzia e la primaria. Ho imparato che il gelato alla nocciola costa di più e che i gelatai ne mettono sempre di meno proprio per questo motivo (o almeno così dicono in giro). Grazie a te sono stato al concerto più bello della mia vita (lascio a voi indovinare quale sia). Avevo sempre immaginato le streghe come donne brutte e cattive; mi hai fatto ricredere: sono solo cattive. Ma ho anche scoperto che sanno mostrare tantissimo affetto e dolcezza quando conta, e che in fondo non sono così antipatiche: è solo il loro modo di dimostrare amore. Spesso mi dai del bambino, ma è grazie a te se sono cresciuto sotto tanti punti di vista: ho sempre voluto fare quello “forte” della coppia, ma in realtà, guardandoti affrontare la vita e superare le difficoltà, mi hai fatto tu da esempio su come si affronta il mondo.
Siamo diversi in tante cose: tu sei più romantica, io devo sforzarmi per esserlo; tu sei creativa, io sono più pratico; tu sei delicata, io sono “materiale”. A me piacciono film e serie TV belle e interessanti, tu guardi solo monnezza; io ascolto cantanti e musica che hanno fatto la storia, tu purtroppo hai bisogno di farti una cultura musicale... e potrei continuare all'infinito. Ma più che differenze, le definirei come il completamento l'uno dell'altra. In realtà abbiamo anche tante cose in comune, ma mi fa più ridere pensare a quelle su cui siamo diversi. Tuttavia, proprio il ridere è una cosa su cui siamo profondamente simili: quando siamo insieme, anche nei momenti più difficili, troviamo sempre qualcosa su cui sorridere. Ti rubo giusto una frase, perché credo riassuma alla perfezione quello che siamo: \textit{“Sempre quello di cui ho bisogno”}.

Ai miei \textbf{zii} e \textbf{cugini}: i primi sono come altri genitori per me, i secondi sono i fratelli e le sorelle che non ho mai avuto. Sono passati gli anni ma ci vogliamo sempre un gran bene, e sarà così fino alla fine. Anche se siamo sparpagliati per l'Italia, non vediamo l'ora di rivederci e di passare del tempo insieme, perché la cosa che ci lega è l'affetto e non il vivere nello stesso posto. Molti di voi mi hanno visto letteralmente crescere e hanno assistito alle tappe più importanti della mia vita, orgogliosi di quello che sono diventato. Magari non lo do a vedere, ma al solo pensiero mi si riempie il cuore di gioia ogni volta. Anche se i nonni non ci sono più, ogni volta che metto piede in casa loro ritorno con la mente a quando ero piccolo e ricordo tutti i bei momenti passati insieme. Un giorno, quando forse vivrò da un'altra parte, so che mi mancherà tanto fare la spesa con \textbf{zia Sofia}, aggiustare il telefono a \textbf{zio Ciccio}, sentire \textbf{Gerardo} e \textbf{Rachele} che predicano e tutto ciò che ha reso la mia infanzia e la mia adolescenza così speciali. Grazie a tutti voi.

A \textbf{Ciccio, Secco, Spritz, Andrea, Giacomo e Deggy, Carlo e Angelo}: in queste ultime settimane mi sono fatto sentire poco, spesso non ho risposto a chiamate e messaggi e probabilmente mi avete trovato più rintronato del solito. Non vi preoccupate, da adesso si torna alla normalità, \textit{“o show insomma”}. È anche grazie a voi e ai momenti passati insieme se non sono impazzito del tutto e sono riuscito a superare questo periodo. Mi avete sopportato quando ero stressato e in paranoia, ricaricandomi a ogni uscita. Non continuerò con i sentimentalismi perché non fanno per noi. Sarete per sempre \textit{“a storj”}.

A \textbf{Martina}: abbiamo affrontato tanti calvari insieme, da Pepper alla tesi, ma mi basta dire una sola cosa per spiegare tutto: \textit{“Mia sorella, di nome e di fatto”}.

A \textbf{Simone, Fabio e Aniello}: vi ringrazio per aver sopportato un brontolone che passava le ore a lamentarsi di tutto ogni volta che passavate per il laboratorio di Ciaramella. Nessuno lo sa, ma in realtà siete quelli che, insieme a me, conoscono il mio progetto di tesi meglio di chiunque altro, dato che ogni volta che mettete piede nel laboratorio ve lo rispiegavo alla lavagna. Mi avete dato un grande aiuto con i vostri consigli e facendomi scoprire nuovi “strumenti” che mi hanno aiutato in questo lavoro. Questa tesi è anche vostra.

A \textbf{Marco, Stefano, Renato e Vincenzo}: con voi (e non solo) abbiamo seguito i corsi, ci siamo fatti paranoie per le lezioni, ci siamo lamentati all'inverosimile e abbiamo bevuto litri di caffè. Siete state le prime vittime dei miei lamenti all'università, ma anche quelli che mi hanno sempre aiutato a superare gli esami dandomi consigli, tranquillizzandomi e, a volte, facendomi da supporto morale. In bocca al lupo per qualsiasi cosa facciate.

A \textbf{Max e Denny}: vi dedico uno spazio separato perché per me la \textbf{“Task Force”} è stata il team di lavoro più efficiente e, allo stesso tempo, bello in cui abbia mai lavorato. Voi eravate due macchine, io invece mi occupavo di rendere il clima più leggero. Anche se non lavoriamo più insieme da anni, l'affetto e la stima che proviamo gli uni per gli altri rimangono intatti. In questi mesi ci siamo sentiti e ci siamo fatti un sacco di risate come ai vecchi tempi. So che forse non lavoreremo più insieme, ma spero di continuare a far uscire Max dal laboratorio in cui si rintana e di andare, un giorno, a trovare Denny in Germania.

Ai ragazzi del gruppo \textbf{“Stas\ae r”}: è da un po' che non ci vediamo, ma vi voglio un gran bene e so che la cosa è reciproca. Si sono riuniti i Co'Sang, si sono riuniti i Club Dogo... penso sia venuto il momento di una reunion anche per noi.

Ai ragazzi del gruppo \textbf{“Ruote di scorta di Cristian”}: ultimamente siamo stati un po' distanti, ognuno preso dai propri impegni, ma quando ci rivediamo torniamo sempre ai tempi del liceo come se non fossimo mai cresciuti. In fondo, nessuno di noi è cambiato davvero. Saremo per sempre la \textbf{5\textsuperscript{a}F}, e per me sarete sempre tra i miei più cari amici.

Al professore \textbf{Maratea}: lo ringrazio per avermi dato l'opportunità di lavorare con lui e per aver dedicato tanto tempo e attenzione al mio lavoro. Un grazie va anche a tutti i professori che ho incontrato durante il mio percorso universitario, perché ognuno di loro mi ha lasciato qualcosa di importante, e all'\textbf{Università Parthenope} in generale.

Infine, riservo una parte per parlare a me stesso. Non voglio ringraziarti per essere arrivato fin qui (in fondo si viene al mondo per vivere al meglio, credo), e non voglio fare troppa filosofia, quindi ti parlerò di musica. Diceva Caparezza: \textit{“Il secondo disco è sempre il più difficile”}. Effettivamente la seconda laurea per te è stata più difficile della prima, e in questi anni ti sei ripetuto più di una volta \textit{“Jodellavitanonhocapitouncazzo”}, come diceva sempre il buon Michele Salvemini. Magari dici così perché non sei in grado di premiarti per i tuoi traguardi, ma la verità è che siamo diventati grandi, siamo ancora io e te, e poteva andare peggio. Non sono parole mie, ma di un altro artista che apprezziamo tanto, quindi ti lascerò con quest'ultima sua barra, che dovrai ricordare per sempre: \textit{“Sii grato”}.